% p102_1; source PDF page 102; manually verified editable table.
\begin{tikzpicture}[x=1cm,y=-1cm,inner sep=0pt,outer sep=0pt]
\path[use as bounding box] (0,0) rectangle (16.000,5.310);
\draw[line width=.4pt] (0.000,0.000) rectangle (1.300,0.905);
\node[align=center,text width=1.040cm,font={\linespread{1}\fontsize{10}{13.50}\selectfont\bfseries}] at (0.650,0.453) {种类};
\draw[line width=.4pt] (1.300,0.000) rectangle (5.100,0.905);
\node[align=center,text width=3.540cm,font={\linespread{1}\fontsize{10}{13.50}\selectfont\bfseries}] at (3.200,0.453) {起因条件};
\draw[line width=.4pt] (5.100,0.000) rectangle (7.900,0.905);
\node[align=center,text width=2.540cm,font={\linespread{1}\fontsize{10}{13.50}\selectfont\bfseries}] at (6.500,0.453) {时间条件};
\draw[line width=.4pt] (7.900,0.000) rectangle (10.600,0.905);
\node[align=center,text width=2.440cm,font={\linespread{1}\fontsize{10}{13.50}\selectfont\bfseries}] at (9.250,0.453) {意思条件};
\draw[line width=.4pt] (10.600,0.000) rectangle (13.300,0.905);
\node[align=center,text width=2.440cm,font={\linespread{1}\fontsize{10}{13.50}\selectfont\bfseries}] at (11.950,0.453) {“不得已”条件};
\draw[line width=.4pt] (13.300,0.000) rectangle (16.000,0.905);
\node[align=center,text width=2.440cm,font={\linespread{1}\fontsize{10}{13.50}\selectfont\bfseries}] at (14.650,0.453) {限度条件};
\draw[line width=.4pt] (0.000,0.905) rectangle (1.300,5.310);
\node[align=left,text width=1.040cm,font={\linespread{1}\fontsize{10}{13.50}\selectfont}] at (0.650,3.108) {正当防卫};
\draw[line width=.4pt] (1.300,0.905) rectangle (5.100,5.310);
\node[align=left,text width=3.540cm,font={\linespread{1}\fontsize{10}{13.50}\selectfont}] at (3.200,3.108) {制止不法侵害。可以为了保护国家法益、公共法益而实施正当防卫};
\draw[line width=.4pt] (5.100,0.905) rectangle (7.900,5.310);
\node[align=left,text width=2.540cm,font={\linespread{1}\fontsize{10}{13.50}\selectfont}] at (6.500,3.108) {不法侵害正在进行，紧迫程度高一些};
\draw[line width=.4pt] (7.900,0.905) rectangle (10.600,5.310);
\node[align=left,text width=2.440cm,font={\linespread{1}\fontsize{10}{13.50}\selectfont}] at (9.250,3.108) {偶然防卫，防卫认识不要说与必要说};
\draw[line width=.4pt] (10.600,0.905) rectangle (13.300,5.310);
\node[align=left,text width=2.440cm,font={\linespread{1}\fontsize{10}{13.50}\selectfont}] at (11.950,3.108) {原则上，没有躲避义务，防卫反击是优先手段，不要求不得已而为之};
\draw[line width=.4pt] (13.300,0.905) rectangle (16.000,5.310);
\node[align=left,text width=2.440cm,font={\linespread{1}\fontsize{10}{13.50}\selectfont}] at (14.650,3.108) {更宽松。防卫明显超过必要限度造成重大损害的，才属于防卫过当};
\end{tikzpicture}
